\subsection{相关工作与现有方案不足}

针对外卖订单通信中的安全问题，学术界和产业界已开展了一定的研究和实践，但现有方案在系统性、安全性和适用性方面仍存在不足。

\textbf{现有外卖平台的隐私保护措施。}以美团、饿了么为代表的外卖平台已推出虚拟号码、隐私面单等隐私保护功能。例如，美团规则中心规定商家和骑手需通过平台生成的虚拟号码联系用户，虚拟号码过期后即无法再行联系\footnote{参见美团规则中心"消费者权益保护"相关条款。}。然而，这些措施主要针对电话通信环节的号码脱敏，并未覆盖订单即时通信的全流程。如前文所述，虚拟号码技术在实际部署中存在短信通道未加密、多次通话后关联真实号码等漏洞；同时，现有平台缺乏对聊天消息内容的加密存储、通信日志的完整性验证以及异常行为的条件溯源机制，难以应对日志篡改、越权访问和纠纷追责等安全需求。

\textbf{通用匿名通信系统。}以Tor（The Onion Router）为代表的匿名通信系统通过多层加密和中间节点转发实现通信匿名性，在隐私保护领域具有重要影响。然而，这类系统的设计目标是通用互联网通信的匿名性，并未考虑外卖订单场景中订单认证、配送业务逻辑和平台治理等特殊需求。在Tor等系统中，通信参与方完全匿名，平台无法在发生纠纷时进行条件溯源，因此不适合直接应用于外卖订单通信场景。

\textbf{学术研究现状。}在学术研究层面，条件隐私保护（Conditional Privacy-Preserving）和可溯源匿名（Traceable Anonymity）已在车联网、电子投票、移动支付等领域得到较多研究。这些研究通常采用群签名、环签名、盲签名或零知识证明等密码学工具实现匿名性与可溯源性之间的平衡。然而，面向外卖订单即时通信场景的条件隐私保护方案研究尚处于起步阶段，现有方案多数停留在理论分析层面，缺少可运行的系统原型和端到端的安全机制闭环。

综上所述，现有方案在外卖订单通信安全方面存在三个主要不足：一是隐私保护措施碎片化，缺少覆盖身份匿名、消息保护和日志审计的系统化设计；二是匿名机制与平台治理之间缺乏平衡，要么过度暴露身份、要么完全不可溯源；三是多数方案缺少工程实现和实际验证。本作品旨在填补这一空白，设计并实现面向外卖订单通信场景的隐私认证与可信审计系统。
